
\documentclass{article}
\usepackage{colortbl}
\usepackage{makecell}
\usepackage{multirow}
\usepackage{supertabular}

\begin{document}

\newcounter{utterance}

\centering \large Interaction Transcript for game `referencegame', experiment `line\_grids\_columns', episode 0 with NEW\_base\_planning\_clp\_chat.
\vspace{24pt}

{ \footnotesize  \setcounter{utterance}{1}
\setlength{\tabcolsep}{0pt}
\begin{supertabular}{c@{$\;$}|p{.15\linewidth}@{}p{.15\linewidth}p{.15\linewidth}p{.15\linewidth}p{.15\linewidth}p{.15\linewidth}}
    \# & $\;$A & \multicolumn{4}{c}{Game Master} & $\;\:$B\\
    \hline

    \theutterance \stepcounter{utterance}  
    & & \multicolumn{4}{p{0.6\linewidth}}{
        \cellcolor[rgb]{0.9,0.9,0.9}{
            \makecell[{{p{\linewidth}}}]{
                \texttt{\tiny{[P1$\langle$GM]}}
                \texttt{You are given three grids, where each of them is 5 by 5 in size.} \\
\texttt{Grids have empty cells marked with "▢" and filled cells marked with "X".} \\
\texttt{Your task is to generate a referring expression that best describes the target grid while distinguishing it from the two other distractor grids.} \\
\texttt{The first grid is the target grid, and the following two grids are the distractors.} \\
\\ 
\texttt{Target grid:} \\
\\ 
\texttt{X □ X X □} \\
\texttt{X □ X X □} \\
\texttt{X □ X X □} \\
\texttt{X □ X X □} \\
\texttt{X □ X X □} \\
\\ 
\texttt{Distractor grid 1:} \\
\\ 
\texttt{□ □ X X X} \\
\texttt{□ □ X X X} \\
\texttt{□ □ X X X} \\
\texttt{□ □ X X X} \\
\texttt{□ □ X X X} \\
\\ 
\texttt{Distractor grid 2:} \\
\\ 
\texttt{X □ □ X X} \\
\texttt{X □ □ X X} \\
\texttt{X □ □ X X} \\
\texttt{X □ □ X X} \\
\texttt{X □ □ X X} \\
\\ 
\texttt{Instruction: Describe the target grid.} \\
\texttt{Generate the referring expression starting with the tag "Expression: " for the given target grid. Omit any other text.} \\
            }
        }
    }
    & \\ \\

    \theutterance \stepcounter{utterance}  
    & \multicolumn{4}{p{0.45\linewidth}}{
        \cellcolor[rgb]{0.95,0.95,0.95}{
            \makecell[{{p{\linewidth}}}]{
                \texttt{\tiny{[P1$|$internal]}}
                \texttt{{[}} \\
\texttt{  \{} \\
\texttt{    "role": "user",} \\
\texttt{    "content": "You are given three grids, where each of them is 5 by 5 in size.\textbackslash{}nGrids have empty cells marked with \textbackslash{}"\textbackslash{}u25a2\textbackslash{}" and filled cells marked with \textbackslash{}"X\textbackslash{}".\textbackslash{}nYour task is to generate a referring expression that best describes the target grid while distinguishing it from the two other distractor grids.\textbackslash{}nThe first grid is the target grid, and the following two grids are the distractors.\textbackslash{}n\textbackslash{}nTarget grid:\textbackslash{}n\textbackslash{}nX \textbackslash{}u25a1 X X \textbackslash{}u25a1\textbackslash{}nX \textbackslash{}u25a1 X X \textbackslash{}u25a1\textbackslash{}nX \textbackslash{}u25a1 X X \textbackslash{}u25a1\textbackslash{}nX \textbackslash{}u25a1 X X \textbackslash{}u25a1\textbackslash{}nX \textbackslash{}u25a1 X X \textbackslash{}u25a1\textbackslash{}n\textbackslash{}nDistractor grid 1:\textbackslash{}n\textbackslash{}n\textbackslash{}u25a1 \textbackslash{}u25a1 X X X\textbackslash{}n\textbackslash{}u25a1 \textbackslash{}u25a1 X X X\textbackslash{}n\textbackslash{}u25a1 \textbackslash{}u25a1 X X X\textbackslash{}n\textbackslash{}u25a1 \textbackslash{}u25a1 X X X\textbackslash{}n\textbackslash{}u25a1 \textbackslash{}u25a1 X X X\textbackslash{}n\textbackslash{}nDistractor grid 2:\textbackslash{}n\textbackslash{}nX \textbackslash{}u25a1 \textbackslash{}u25a1 X X\textbackslash{}nX \textbackslash{}u25a1 \textbackslash{}u25a1 X X\textbackslash{}nX \textbackslash{}u25a1 \textbackslash{}u25a1 X X\textbackslash{}nX \textbackslash{}u25a1 \textbackslash{}u25a1 X X\textbackslash{}nX \textbackslash{}u25a1 \textbackslash{}u25a1 X X\textbackslash{}n\textbackslash{}nInstruction: Describe the target grid.\textbackslash{}nGenerate the referring expression starting with the tag \textbackslash{}"Expression: \textbackslash{}" for the given target grid. Omit any other text."} \\
\texttt{  \},} \\
\texttt{  \{} \\
\texttt{    "role": "user",} \\
\texttt{    "content": "You are a skilled planner. Above you see the game history. What would be the best next action? Openly reason about the next step, then produce your next action exactly as the game expects it preceded by the separator 'ACT:' on its own line. If multiple actions are required, write ACT: first and then all actions."} \\
\texttt{  \}} \\
\texttt{{]}} \\
            }
        }
    }
    & & \\ \\

    \theutterance \stepcounter{utterance}  
    & \multicolumn{4}{p{0.45\linewidth}}{
        \cellcolor[rgb]{0.95,0.95,0.95}{
            \makecell[{{p{\linewidth}}}]{
                \texttt{\tiny{[P1$|$internal]}}
                \texttt{The target grid has a distinct pattern: the first column is entirely filled with "X", the second column is entirely "□", and the third, fourth, and fifth columns are "X", "X", "□" respectively, repeating across all rows.} \\
\\ 
\texttt{In contrast:} \\
\texttt{{-} Distractor 1 has the first two columns as "□", and the rest as "X".} \\
\texttt{{-} Distractor 2 has the first column as "X", second as "□", third as "□", and the last two as "X".} \\
\\ 
\texttt{The key distinguishing feature of the target grid is that the second column is all "□" and the first column is all "X", with the rest of the pattern being consistent across rows.} \\
\\ 
\texttt{Expression: The grid where the first column is all X, the second column is all □, and the third, fourth, and fifth columns are X, X, □ respectively.} \\
            }
        }
    }
    & & \\ \\

    \theutterance \stepcounter{utterance}  
    & \multicolumn{4}{p{0.6\linewidth}}{
        \cellcolor[rgb]{0.9,0.9,0.9}{
            \makecell[{{p{\linewidth}}}]{
                \texttt{\tiny{[P1$\rangle$GM]}}
                \texttt{The target grid has a distinct pattern: the first column is entirely filled with "X", the second column is entirely "□", and the third, fourth, and fifth columns are "X", "X", "□" respectively, repeating across all rows.} \\
\\ 
\texttt{In contrast:} \\
\texttt{{-} Distractor 1 has the first two columns as "□", and the rest as "X".} \\
\texttt{{-} Distractor 2 has the first column as "X", second as "□", third as "□", and the last two as "X".} \\
\\ 
\texttt{The key distinguishing feature of the target grid is that the second column is all "□" and the first column is all "X", with the rest of the pattern being consistent across rows.} \\
\\ 
\texttt{Expression: The grid where the first column is all X, the second column is all □, and the third, fourth, and fifth columns are X, X, □ respectively.} \\
            }
        }
    }
    & & \\ \\

    \theutterance \stepcounter{utterance}  
    & & & \multicolumn{2}{p{0.3\linewidth}}{
        \cellcolor[rgb]{0.9,0.9,0.9}{
            \makecell[{{p{\linewidth}}}]{
                \texttt{\tiny{[GM$|$GM]}}
                \texttt{Invalid generated expression} \\
            }
        }
    }
    & & \\ \\

\end{supertabular}
}

\end{document}
